\documentclass[10pt]{article}
% Math Packages
\usepackage{amsmath, mathtools}
\usepackage{amssymb}
\usepackage{amsthm}
\usepackage{amsfonts}
\usepackage{bbm}
\usepackage{breqn}
\usepackage[margin=1in]{geometry}
\usepackage{graphicx}
\usepackage{tikz}
\usetikzlibrary{arrows.meta}
\usetikzlibrary{calc}
\usepackage{forest}
\usepackage{tikz-qtree}
\graphicspath{ {./images/} }
\usepackage{hyperref}
\usepackage[capitalize]{cleveref}
\usepackage[shortlabels]{enumitem}
\usetikzlibrary{arrows,matrix,positioning}
\usepackage{multicol}

\usepackage{listings}
\usepackage{xcolor}

\crefname{problem}{Problem}{Problems}
\Crefname{problem}{Problem}{Problems}

\lstset{
  language=R,
  basicstyle=\ttfamily\footnotesize,
  keywordstyle=\color{blue},
  commentstyle=\color{gray},
  stringstyle=\color{red!70!black},
  framesep=2pt,
  framexleftmargin=0pt, 
  frame=single,
  breaklines=false,
  showstringspaces=false
}
% for the pipe symbol
\usepackage[T1]{fontenc}

% Colorful Notes
\usepackage{color} \definecolor{Red}{rgb}{1,0,0} \definecolor{Blue}{rgb}{0,0,1}
\definecolor{Purple}{rgb}{.5,0,.5} \def\red{\color{Red}} \def\blue{\color{Blue}}
\def\gray{\color{gray}} \def\purple{\color{Purple}}
\newcommand{\rnote}[1]{{\red [#1]}} % \rnote{foo} gives '[foo]' in red
\newcommand{\pnote}[1]{{\purple [#1]}} % \pnote{foo} gives '[foo]' in purple
\newcommand{\bnote}[1]{{\blue #1}} % \bnote{foo} gives 'foo' in blue
\newcommand{\gnote}[1]{{\gray #1}} % \gnote{foo} gives 'foo' in gray
\newcommand{\Max}[1]{{\purple [#1]}} % \bnote{foo} then 'foo' is blue

% Claim numbering (the counter restarts after each proof environment)
\newcounter{claimcount}
\setcounter{claimcount}{0}
\newenvironment{claim}{\refstepcounter{claimcount}\par\addvspace{\medskipamount}\noindent\textbf{Claim \arabic{claimcount}:}}{}
\usepackage{etoolbox}
\AtBeginEnvironment{proof}{\setcounter{claimcount}{0}}
\newenvironment{claimproof}{\par\addvspace{\medskipamount}\noindent\textit{Proof of Claim  \arabic{claimcount}.}}{\hfill\ensuremath{\qedsymbol} \tiny{Claim}

  \medskip}
% Add claim support to cleverref
\crefname{claimcount}{Claim}{Claims}

% Math Environments
\newtheorem{theorem}{Theorem}
\newtheorem{assumption}[theorem]{Assumption}
\newtheorem{lemma}[theorem]{Lemma}
\newtheorem{proposition}[theorem]{Proposition}
\newtheorem{corollary}[theorem]{Corollary}
\newtheorem{question}[theorem]{Question}
\theoremstyle{definition}
\newtheorem*{definition}{Definition}
\newtheorem{remark}[theorem]{Remark}
\newtheorem{example}[theorem]{Example}
\newtheorem{notation}[theorem]{Notation}
\newtheorem{problem}[theorem]{Problem}

% Matrices and Column Vectors. 
\usepackage{stackengine}
\setstackgap{L}{1.0\normalbaselineskip}
\usepackage{tabstackengine}
\setstackEOL{;}% row separator
\setstackTAB{,}% column separator
\setstacktabbedgap{1ex}% inter-column gap
\setstackgap{L}{1.5\normalbaselineskip}% inter-row baselineskip
\let\nmatrix\bracketMatrixstack  %Usage: \nmatrix{1,2,3\4,5,6}
\newcommand\cv[1]{\setstackEOL{,}\bracketMatrixstack{#1}} %usage: \cv{1,2,3}

% Custom Math Coqmmands
\newcommand{\vt}{\vskip 5mm} % vertical space
\newcommand{\fl}{\noindent\textbf} % first line
\newcommand{\Fl}{\vt\noindent\textbf} % first line with space above
\newcommand{\norm}[1]{\left\lVert#1\right\rVert} % norm
\newcommand{\pnorm}[1]{\left\lVert#1\right\rVert_p} % p-norm
\newcommand{\qnorm}[1]{\left\lVert#1\right\rVert_q} % q-norm
\newcommand{\1}[1]{\textbf{1}_{\left[#1\right]}} % indicator function 
\def\limn{\lim_{n\to\infty}} % shortcut for lim as n-> infinity
\def\sumn{\sum_{n=1}^{\infty}} % shortcut for sum from n=1 to infinity
\def\sumkn{\sum_{k=1}^{n}} % shortcut for sum from k=1 to n
\def\sumin{\sum_{i=1}^{n}} % shortcut for sum from i=1 to n
\def\SAs{\sigma\text{-algebras}} % shortcut for $\sigma$-algebras
\def\SA{\sigma\text{-algebra}} % shortcut for $\sigma$-algebra
\def\Ft{\mathcal{F}_t} % time-indexed sigma-algebra (t)
\def\Fs{\mathcal{F}_s} % time-indexed sigma-algebra (s)
\def\F{\mathcal{F}} % sigma-algebra
\def\G{\mathcal{G}} % sigma-algebra
\def\R{\mathbb{R}} % Real numbers
\def\N{\mathbb{N}} % Natural numbers
\def\Z{\mathbb{Z}} % Integers
\def\E{\mathbb{E}} % Expectation
\def\P{\mathbb{P}} % Probability
\def\Q{\mathbb{Q}} % Q probability
\def\dist{\text{dist}} %Text 'dist' for things like 'dist(x,y)'
\newcommand{\indep}{\perp \!\!\! \perp}  %independence symbol
\def\Var{\mathrm{Var}} % Variance
\def\tr{\mathrm{tr}} % trace

% Brackets and Parentheses
\def\[{\left [}
    \def\]{\right ]}
% \def\({\left (}
%   \def\){\right )}

\usepackage{color}
\definecolor{Red}{rgb}{1,0,0}
\definecolor{Blue}{rgb}{0,0,1}
\definecolor{Purple}{rgb}{.5,0,.5}
\def\red{\color{Red}}
\def\blue{\color{Blue}}
\def\gray{\color{gray}}
\def\purple{\color{Purple}}
\definecolor{RoyalBlue}{cmyk}{1, 0.50, 0, 0}
\newcommand{\dempfcolor}[1]{{\color{RoyalBlue}#1}} 
\newcommand{\demph}[1]{\textcolor{RoyalBlue}{\textbf{\slshape #1}}} % Slanted

\usepackage{emoji}
% comment exactly one of the following line to show / hide the solutions
% \newcommand{\solution}[1]{{\purple #1}} % uncomment to show the solutions
\newcommand{\solution}[1]{} % uncomment to hide the solutions




\begin{document}
\begin{center}
  \section*{Math 472: Homework 11}
  \textit{Due Wednesday, May 6}
\end{center}

\begin{problem}[Exercise 10.117]
  Lord Rayleigh was one of the earliest scientists to study the density of
  nitrogen.\footnote{Lord Rayleigh, On an Anomaly Encountered in
    Determinations of the Density of Nitrogen Gas, \textit{Proceedings of the
      Royal Society of London}, (1894),
    \url{http://www.jstor.org/stable/115479}} He employed two types of methods
  to produce nitrogen gas: (1) a process to derive nitrogen from atmospheric
  gas via removal of oxygen, and (2) a process which produced nitrogen gas
  from chemical compounds not involving atmospheric gas.

  In his studies, he noticed something peculiar. The nitrogen densities
  produced from chemical compounds tended to be smaller than the densities of
  nitrogen produced from atmospheric gas. Lord Rayleigh's measurements are
  given in the following table. These measurements correspond to the mass of
  nitrogen filling a flask of specified volume under specified temperature and
  pressure.
  \begin{center}
    \centering
    \begin{tabular}{cc}
      \hline
      \textbf{Compound Chemical} & \textbf{Atmosphere} \\
      \hline
      2.30143 & 2.31017 \\
      2.29890 & 2.30986 \\
      2.29816 & 2.31010 \\
      2.30182 & 2.31001 \\
      2.29869 & 2.31024 \\
      2.29940 & 2.31010 \\
      2.29849 & 2.31028 \\
      2.29889 & 2.31163 \\
      2.30074 & 2.30956 \\
      2.30054 &         \\
      \hline
    \end{tabular} 
  \end{center}
  \begin{enumerate}[(a)]
    \item For the measurements from the chemical compound, the sample mean is
    $\overline{y} = 2.29971$ and the standard error $s = .001310$; for the measurements from the
    atmosphere, $\overline{y} = 2.310217$ and $s = .000574$. Is there
    sufficient evidence to indicate a difference in the mean mass of nitrogen
    per flask for chemical compounds and air? What can be said about the
    p-value associated with your test?
    \item Find a 95\% confidence interval for the difference in mean mass of
    nitrogen per flask for chemical compounds and air.
    \item Based on your answer to part (b), at the $\alpha = .05$ level of
    significance, is there sufficient evidence to indicate a difference in mean
    mass of nitrogen per flask for measurements from chemical compounds and
    air?
    \item Is there any conflict between your conclusions in parts (a) and (b)?
    Although the difference in these mean nitrogen masses is small, Lord
    Rayleigh emphasized this difference rather than ignoring it, and this led
    to the discovery of inert gases like Argon in the atmosphere.
  \end{enumerate}
\end{problem}


\begin{problem}
  The octane number $Y$ of refined petroleum is related to the temperature $x$ of
  the refining process, but it is also related to the particle size of the
  catalyst. An experiment with a small-particle catalyst gave a fitted
  least-squares line of
  \begin{equation*}
    \widehat{y} = 9.360 + .155x,
  \end{equation*}
  with $n=31$, $\Var(\widehat{\beta}_{1})= (.0202)^{2}$, and $\mathrm{SSE} =
  2.04$. An independent experiment with a large-particle catalyst gave
  \begin{equation*}
    \widehat{y}= 4.265 +.190x,
  \end{equation*}
  with $n=11$, $\Var(\widehat{\beta}_{1})= (.0193)^{2}$ and $\mathrm{SSE}
  =1.86$.

  Test the hypotheses that the slopes are significantly different from
  zero, with each test at the significance level of $.05$.
\end{problem}
% \begin{problem}[Exercise 10.97]
%   Let $Y_{1},\ldots,Y_{n}$ be independent and identically distributed random
%   variables with probability mass function $p(y\mid \theta)$ given by
%   \begin{equation*}
%     p(1\mid \theta)  = \theta^{2},
%     \quad
%     p(2\mid \theta)  = 2\theta(1-\theta)
%     \quad \text{and} \quad
%     p(3\mid \theta)  = (1-\theta)^{2},
%   \end{equation*}
%   where $\theta\in (0,1)$. Let $N_{i}$ denote the number of observations equal
%   to $i$ for $i=1,2,3$.
%   \begin{enumerate}[(a)]
%     \item Derive the likelihood function $L(\theta)$ as a function of
%     $N_{1},N_{2},$ and $N_{3}$.
%     \item Find the most powerful test for testing $H_{0}:\theta=\theta_{0}$
%     versus $H_{a}:\theta=\theta_{a}$, where $\alpha_{a}>\alpha_{0}$. Show that
%     your test specifies that $H_{0}$ be rejected for certain values of
%     $2N_{1}+N_{2}.$
%     \item How do you determine the value of $k$ so that the test has nominal
%     level $\alpha$? You do not need to do the actual computation. A clear
%     description of how to determine $k$ is adequate.
%     \item Is the test derived in parts (a)-(c) uniformly most powerful for
%     testing $H_{0}:\theta=\theta_{0}$ versus $H_{a}:\theta>\theta_{0}$? Why or
%     why not?
%   \end{enumerate}
% \end{problem}

\begin{problem}
  A recent study of collected fossil data\footnote{P. Monarrez, et al. Reduced
    strength and increased variability of extinction selectivity during mass
    extinctions, \textit{R Soc Open Sci} (2023)
    \url{https://royalsocietypublishing.org/rsos/article/10/9/230795/92261}}
  investigated the relationship between species' traits and their risk of
  going extinct. In this problem we'll analyze a subset of their dataset,
  specifically a subset of 550 trilobita species, with fossil samples taken
  from various epochs. In particular, we'll look at data consisting of
  $n_{1}=40$ samples of trilobite fossils which date back to epochs of known
  mass extinctions, and compare them with $n_{2}=13070$ samples from
  non-mass-extinction epochs.

  Download the trilobita fossil datafiles from

  \begin{quote}
    {\footnotesize\url{https://github.com/max-hill/math-472/blob/main/datasets/trilobita-fossil-mass-extinctions.csv}}

    {\footnotesize\url{https://github.com/max-hill/math-472/blob/main/datasets/trilobita-fossil-background.csv}}

  \end{quote}
  The first dataset consists of fossil samples dating to geologic epochs
  coinciding with mass extinction events. The second dataset consists of
  samples from ``ordinary times'' (non-mass-extinction epochs).
  
  Each row of these datasets correspond to a fossil sample. The column
  labelled \texttt{extinct} is a boolean variable indicating whether the
  species represented by that fossil went extinction during that epoch. The
  columns \texttt{log\_range} and \texttt{log\_bodysize} correspond to the
  relative geographic range of the species and the body size of the fossil
  sample, respectively. (Here, \texttt{extinct} is the response variable $Y$;
  the covariates $x_{1}$ and $x_{2}$ are \texttt{log\_range} and
  \texttt{log\_bodysize}).
  
  
  \begin{enumerate}[(a)]
    \item Load the background dataset as a dataframe in \texttt{R} and run a
    logistic regression using

\begin{verbatim}
  df_BG = read.csv("trilobita-fossil-background.csv")
  fit_BG <- glm(extinct ~ log_range + log_bodysize, data=df_BG, family=binomial())
\end{verbatim}

    \noindent You can view the results of the logistic regression by running:

\begin{verbatim}
  fit_BG
  summary(fit_BG)
\end{verbatim}

    Interpret your results. What is the best fit logistic regression model
    that you got? What were the estimated coefficients. Which estimates were
    statistically significant and which were not? How do you interpret the
    $p$-values?

    \item Can you draw any inferences about the relationship between (1) body
    size and extinction risk, or (2) between geographic range and body size
    for this dataset? Consider the signs and magnitudes of your estimates
    $\widehat{\beta}_{0},\widehat{\beta}_{1},\widehat{\beta}_{2}$ and their
    $p$-values.

    \item Perform a logistic regression analysis of the mass-extinction epoch
    dataset. Answer questions (a) and (b) for the fossil samples from the mass
    extinction epochs.

    \item Make a comparison between the two logistic regression analyses. How
    does the effect of body size and geographic range on extinction risk
    differ between (1) eras of mass extinction and (2) eras without mass
    extinction? Can you draw any biologically interesting conclusions? Explain
    your results in a way that a biologist would understand.
  \end{enumerate}
\end{problem}

\end{document}
% How your world is turning faster, so much higher you can go